\section{Cutoffs for primitive holomorphic newforms}
\label{hnew:section}

The angular Fourier transform gives the corresponding cutoff for every
fixed primitive holomorphic cuspidal newform.  Weight one then gives
a cutoff for every odd irreducible two-dimensional Artin representation
over $\mathbb Q$ through the classical modularity theorem.

Fix a normalized primitive holomorphic cuspidal newform
$f\in S_k(\Gamma_0(N),\xi)$, where $k\ge1$ is an integer,
$N\ge1$, and $\xi$ is the nebentypus, extended by zero at integers
not coprime to $N$.  Write
\[
 f(z)=\sum_{n\ge1}\lambda_f(n)n^{(k-1)/2}e^{2\pi inz},
 \qquad \lambda_f(1)=1,\qquad
 a=\frac{k-1}{2},\quad \nu=k-1,\quad c=\frac{2\pi}{\sqrt N}.
\]
The coefficients may be complex.  Denote by $\bar f$ the primitive
dual form, with $\lambda_{\bar f}(n)=\overline{\lambda_f(n)}$,
the same weight and level, and nebentypus $\bar\xi$.
All constants below may depend on the fixed form and, where applicable,
the fixed concentration threshold; no uniformity
in its weight, level, or nebentypus is asserted.  Put
\begin{equation}\label{hnew:completed}
 L_f(s)=\sum_{n\ge1}\lambda_f(n)n^{-s},\qquad
 \Lambda_f(s)=c^{-s}\Gamma(s+a)L_f(s).
\end{equation}
Both $L_f$ and the completed function are entire, and
\[
 \Lambda_f(s)=\varepsilon_f\Lambda_{\bar f}(1-s),\qquad
 |\varepsilon_f|=1,\qquad
 \varepsilon_{\bar f}=\overline{\varepsilon_f}.
\]
The standard completed normalization
$N^{s/2}\Gamma_{\mathbb C}(s+a)L_f(s)$ differs from
$\Lambda_f(s)$ by the constant $2(2\pi)^{-a}$, which cancels
from the functional equation; see
\citet[pp.~1--3 of the author preprint]{booker2019}.

On the raw space $\mathcal H=L^2((0,\infty),du)$ define
\[
 (\mathcal H_{c,\nu}\psi)(v)
   =c\int_0^\infty J_\nu(2c\sqrt{uv})\psi(u)\,du,\qquad
 (U\psi)(t)=e^{t/2}\psi(e^t),\qquad H_D^s=U^{-1}H^s(\mathbb R).
\]
The transform is a unitary involution, as verified below.
Set
\[
 P_\Lambda=1_{(0,\Lambda)},\quad
 B_\Lambda=\mathcal H_{c,\nu}P_\Lambda\mathcal H_{c,\nu},\quad
 A_\Lambda=P_\Lambda B_\Lambda P_\Lambda,\quad L=\log\Lambda.
\]
For a fixed $0<\vartheta<1$ and for the auxiliary deep cutoff put
\begin{equation}\label{hnew:selections}
 \Pi_\Lambda^{(\vartheta)}=1_{[\vartheta,1]}(A_\Lambda),
 \qquad
 \Pi_\Lambda^{\mathrm d}=1_{[1-\Lambda^{-56},1]}(A_\Lambda).
\end{equation}
Define the theta operator first on smooth compact profiles by
\begin{equation}\label{hnew:theta}
 (E_f\psi)(t)=e^{t/2}\sum_{n\ge1}\lambda_f(n)\psi(ne^t).
\end{equation}
Its closed multiplier realization on the graph domain used below is
$\widehat{E_f\psi}(\tau)=L_f(1/2+i\tau)\widehat{U\psi}(\tau)$.
Here $\widehat g(\tau)=\int_{\mathbb R}g(t)e^{i\tau t}dt$,
with Plancherel measure $d\tau/(2\pi)$.

For either selection let $j_\Lambda:\mathbb C^{r_\Lambda}\to\mathcal H$
be any raw orthonormal column onto its range.  The full image Gram
matrix, orthonormal theta column, and image projection are
\begin{equation}\label{hnew:gram-definition}
 G_\Lambda=(E_fj_\Lambda)^*E_fj_\Lambda,\qquad
 Z_\Lambda=E_fj_\Lambda G_\Lambda^{-1/2},\qquad
 Q_\Lambda=Z_\Lambda Z_\Lambda^*.
\end{equation}
Their existence and independence of the raw basis follow from the
Gram bound below.  There is no principal sector or polar-coordinate
completion, including when $k=1$ and $\nu=0$.
In particular neither $\psi(0)=0$ nor $\int\psi=0$ is imposed.
Put $S_L=1_{[-L,L]}$ and
$D_\Lambda=S_L-S_LQ_\Lambda S_L$.  For an orthonormal image basis
$z_1,\ldots,z_{r_\Lambda}$ define the positive measure
\begin{equation}\label{hnew:measure}
 d\mu_\Lambda(\tau)=
 \left(2L-\sum_i|\widehat{S_Lz_i}(\tau)|^2\right)\frac{d\tau}{2\pi}.
\end{equation}
Its positivity is Bessel's inequality.

For a nontrivial zero $\rho=\beta+i\gamma$ of $L_f$, with its
analytic multiplicity $m_\rho$, let $\mu_\rho$ be the Cauchy
probability measure
\[
 d\mu_\rho(\tau)=
 \frac{|\beta-1/2|\,d\tau}
 {\pi((\tau-\gamma)^2+|\beta-1/2|^2)}
 \quad(\beta\ne1/2),\qquad
 \mu_\rho=\delta_\gamma\quad(\beta=1/2).
\]
Write $\mathfrak b_f=\sum_\rho m_\rho\mu_\rho$ and
$\mathfrak B_f(h)=\int H(i\tau)d\mathfrak b_f(\tau)$, where
$H(z)=\int_{\mathbb R}h(t)e^{zt}dt$.

\begin{theorem}[The holomorphic newform cutoff]
\label{hnew:theorem}
For every fixed primitive form $f$ as above and every fixed
$0<\vartheta<1$, the full-Gram cutoff
$Q_\Lambda^{(\vartheta)}$ has the following properties.
\begin{enumerate}
\item Its raw selected rank is
$4\Lambda/\sqrt N+O_{f,\vartheta}(\log(2+\Lambda))$, and
$G_\Lambda^{-1}=O_{f,\vartheta}(\Lambda^2\log^2(2+\Lambda))$.
\item The measures \eqref{hnew:measure} converge to $\mathfrak b_f$
against every Schwartz frequency test.  For some fixed $C_{f,\vartheta},A_f>0$,
\[
 \mu_\Lambda([T,T+1])\le C_{f,\vartheta}(2+|T|)^{A_f}
       \qquad(\Lambda\ge2,\ T\in\mathbb R).
\]
\item If $(V(h)g)(t)=\int_{\mathbb R}h(r)g(t-r)dr$, then for
every $h\in C_c^\infty(\mathbb R)$,
\begin{equation}\label{hnew:trace-theorem}
 2\log\Lambda\,h(0)-\operatorname{Tr}(Q_\Lambda V(h))
             \longrightarrow\mathfrak B_f(h).
\end{equation}
\end{enumerate}
The measure and trace assertions also hold for the deep selection
in \eqref{hnew:selections}.  No hypothesis on the positions,
multiplicities, or separation of the nontrivial zeros is imposed.
If $L_f$ is zero free in $\Re s>\sigma_*$, with
$1/2\le\sigma_*<1$ and $\delta=\sigma_*-1/2$, then
\[
 |\mathfrak B_f(h)-\mathfrak W_f(h)|
 \le2\delta\mathcal N_\delta(h)
           \sum_\rho\frac{m_\rho}{1+\gamma^2},\qquad
 \mathfrak W_f(h)=\sum_\rho m_\rho H(\rho-1/2),
\]
where
\[
 \mathcal N_\delta(h)=|h(0)|+
 \int_{\mathbb R}e^{\delta|t|}
 \bigl(|t||h''(t)|+2|h'(t)|+(1+\delta^2)|t||h(t)|\bigr)dt.
\]
\end{theorem}

The leading term is $2\log\Lambda$, despite the degree two
$L$-function and rank $4\Lambda/\sqrt N+O(\log\Lambda)$: the logarithmic
window has length $2\log\Lambda$ and one scalar theta channel.
There are no terms $H(1/2)$ or $H(-1/2)$.
We prove the theorem after establishing the required inputs.

\subsection{The angular transform and its sharp cutoff estimates}

\begin{lemma}[Normalization and domain]\label{hnew:transform}
The transform $\mathcal H_{c,\nu}$ is a unitary involution.  With
$\mathcal M\psi(s)=\int_0^\infty\psi(u)u^{s-1}du$, one has
\begin{equation}\label{hnew:mellin-hankel}
 \mathcal M(\mathcal H_{c,\nu}\psi)(s)
  =g_{c,\nu}(s)\mathcal M\psi(1-s),\qquad
 g_{c,\nu}(s)=c^{1-2s}
      \frac{\Gamma(s+a)}{\Gamma(1-s+a)}.
\end{equation}
Consequently
\begin{equation}\label{hnew:poisson}
 E_f\mathcal H_{c,\nu}=\varepsilon_f J E_{\bar f},\qquad
 (Jg)(t)=g(-t).
\end{equation}
The transform reverses raw dilations and preserves every $H_D^s$.
Its smooth angular core consists of profiles $u^ab(u)$ near zero,
with $b$ smooth, and with Schwartz decay at infinity.
\end{lemma}
\begin{proof}
The lift $\psi(u)\mapsto e^{i\nu\theta}\psi(r^2)$ has norm
$\sqrt\pi$ times the raw norm and identifies the raw space with
the angular-frequency-$\nu$ subspace of $L^2(\mathbb R^2)$.
Euclidean Fourier transformation preserves that subspace.  Removing
its constant phase $(-i)^\nu$ and making the fixed radial rescaling
gives $\mathcal H_{c,\nu}$ and proves unitarity and involutivity.
On a convergence strip, substitute $z=2c\sqrt{uv}$ and use
\[
 \int_0^\infty z^{2s-1}J_\nu(z)\,dz
     =2^{2s-1}\frac{\Gamma(s+\nu/2)}{\Gamma(1-s+\nu/2)}
\]
to obtain \eqref{hnew:mellin-hankel}.  Continuation and the unitary
Mellin realization give it on the critical line.
Equation \eqref{hnew:completed} gives
$L_f(s)g_{c,\nu}(s)=\varepsilon_fL_{\bar f}(1-s)$.
Mellin inversion proves \eqref{hnew:poisson} on the angular core:
its Mellin transforms decrease rapidly on closed strips to the
right of $-a$, permitting a contour shift from $\Re s>1$ to
the critical line.  A shift slightly to the left also proves
square integrability at the lower end.  The graph bound in
Lemma~\ref{hnew:arithmetic} extends the identity to the domains used.
Finally, for $(D_r\psi)(u)=r^{-1/2}\psi(u/r)$, substitution gives
$\mathcal H_{c,\nu}D_r=D_{1/r}\mathcal H_{c,\nu}$.
The dilation generator is therefore reversed, proving the Sobolev
assertion.  Smoothness of the angular lift at the origin gives
the stated form of its core.
\end{proof}
\begin{proposition}[Fixed-order geometric estimates]\label{hnew:geometry}
For fixed $\nu$ and $c$,
\begin{align}
 \operatorname{Tr}A_\Lambda
    &=\frac{2c\Lambda}{\pi}+O_{c,\nu}(\log(2+\Lambda)),\label{hnew:trace}\\
 \operatorname{Tr}(A_\Lambda-A_\Lambda^2)
    &=O_{c,\nu}(\log(2+\Lambda)).\label{hnew:variance}
\end{align}
Every fixed-threshold selected vector has a band extension
\[
 W_\Lambda\psi
 =B_\Lambda A_\Lambda^{-1}\psi
 =\mathcal H_{c,\nu}P_\Lambda F_\Lambda^{-1}\psi,\qquad
 F_\Lambda=P_\Lambda\mathcal H_{c,\nu}P_\Lambda,
\]
such that $P_\Lambda W_\Lambda\psi=\psi$ and
$\|W_\Lambda\|\le\vartheta^{-1/2}$.  For some fixed $d>0$,
\begin{equation}\label{hnew:raw-mass-graph}
 \int_{d/\Lambda}^\Lambda|\psi(u)|^2du
   \ge\tfrac12\|\psi\|_2^2,\qquad
 \|\psi\|_{H_D^s}\le C_s\Lambda^s\|\psi\|_2
       \quad(0<s<1/2).
\end{equation}
For a smooth strict-margin $p$, equal to one near zero and supported
in $[0,b]$, $b<1$, write
$M_\Lambda=p(u/\Lambda)$,
$B_\Lambda^p=\mathcal H_{c,\nu}M_\Lambda\mathcal H_{c,\nu}$, and
$K_\Lambda^p=M_\Lambda B_\Lambda^pM_\Lambda$.  For either selection,
every $0\le\sigma<1/2$, and every $R$,
\begin{equation}\label{hnew:interior}
 \|(1-\Pi_\Lambda)K_\Lambda^p\|_{L^2\to H_D^\sigma}
                  =O_{p,\sigma,R}(\Lambda^{-R}).
\end{equation}
For the deep selection and $0\le\sigma<s<1/2$,
\begin{equation}\label{hnew:deep-raw-tail}
 \|(1-P_\Lambda)W_\Lambda\psi\|_{H_D^\sigma}
 \le C\Lambda^{\sigma-28(1-\sigma/s)}\|\psi\|_2.
\end{equation}
\end{proposition}
\begin{proof}
The change $r=2c\sqrt{\Lambda u}$ identifies the band kernel with
\[
 K_\nu(r,s)=\sqrt{rs}\int_0^1vJ_\nu(rv)J_\nu(sv)dv.
\]
Lommel's identity and its diagonal limit give
\[
 K_\nu(r,s)=\frac{\sqrt{rs}
 [sJ_{\nu+1}(s)J_\nu(r)-rJ_{\nu+1}(r)J_\nu(s)]}{s^2-r^2},
 \quad
 K_\nu(r,r)=\frac r2[J_\nu(r)^2-J_{\nu-1}(r)J_{\nu+1}(r)].
\]
The fixed-order asymptotic gives $K_\nu(r,r)=\pi^{-1}+O(r^{-1})$
for $r\ge1$.  The series at zero gives $K_\nu(r,r)=O(r^{2\nu+1})$.
The off-diagonal formula and Cauchy--Schwarz imply
\[
 |K_\nu(r,s)|\le C\min(1,|r-s|^{-1})\quad(r,s\ge1),\qquad
 |K_\nu(r,s)|\le Cr^{\nu+1/2}s^{-1}\quad(r<1,\ s\ge2).
\]
Integrating the diagonal to $Z=2c\Lambda$ proves \eqref{hnew:trace}.
Integrating $|K_\nu(r,s)|^2$ over $r<Z<s$ proves
\eqref{hnew:variance}; the large-variable integral is
$O(\int_0^Z(1+Z-r)^{-1}dr)$.
The scalar inequality
$|1_{x\ge\vartheta}-x|\le C_\vartheta x(1-x)$ proves the
fixed-threshold rank assertion.

Since $F_\Lambda^2=A_\Lambda$ and $F_\Lambda$ commutes with its
spectral selections, the two displayed formulas for $W_\Lambda$
agree.  Moreover
$\|W_\Lambda\psi\|_2^2=\langle\psi,A_\Lambda^{-1}\psi\rangle$.
For a band vector $w$, Cauchy--Schwarz in the Hankel integral gives
$|w(u)|\le C\sqrt\Lambda\|w\|_2$.  This proves the mass assertion
by choosing $d$ sufficiently small.  Termwise integration of the
Bessel series gives $w(u)=u^ab(u)$ with $b$ analytic near zero.
In particular $w'(u)=O(u^{a-1})$ when $a>0$.
When $a=0$, $w=b$ is analytic and $w'$ is bounded near zero.
For positive half-integral $a$, bounded ordinary derivatives of
every order at zero are neither claimed nor needed.

For the angular lift, the Fourier band radius is
$(c/\pi)\sqrt\Lambda$.  Plancherel therefore bounds its gradient
by $2c\sqrt\Lambda$ times its norm.  Since $u\partial_u$
corresponds to one half the radial dilation derivative, its norm
on $u<\Lambda$ is $O(\Lambda)\|w\|_2$.
The logarithmic profile tends to zero at its lower end and has
one jump at $L$.  The trace inequality
$|U\psi(L-)|^2\le2\|U\psi\|_2\|(U\psi)'\|_2$ and the Fourier
formula for that zero extension give \eqref{hnew:raw-mass-graph}
exactly for $s<1/2$.

The soft multipliers lift to multiplication by $p(|z|^2/\Lambda)$
and its Fourier counterpart; both preserve the angular sector.
The fixed gap between the soft and hard Fourier disks, and
Schwartz decay of the Fourier transform of the first multiplier,
give a Schur bound $O_R(\Lambda^{-R})$ on the hard band tail.
For a time-supported vector $v$,
\[
 \|(1-\Pi_\Lambda)v\|_2^2
 \le\begin{cases}
 (1-\vartheta)^{-1}\|(1-B_\Lambda)v\|_2^2,
       &\Pi_\Lambda=\Pi_\Lambda^{(\vartheta)},\\
 \Lambda^{56}\|(1-B_\Lambda)v\|_2^2,
       &\Pi_\Lambda=\Pi_\Lambda^{\mathrm d}.
 \end{cases}
\]
The arbitrary power absorbs this fixed deep loss.  Interpolating
with \eqref{hnew:raw-mass-graph} proves \eqref{hnew:interior}.
This is the precise order-$\nu$ replacement for
Propositions~\ref{prop:hankel-interior-approximation}
and~\ref{prop:hankel-high-interior}.

For the deep extension,
$\|(1-P_\Lambda)W_\Lambda\psi\|_2^2
\le\Lambda^{-56}(1-\Lambda^{-56})^{-1}\|\psi\|_2^2$.
The identity involving $F_\Lambda^{-1}$, dilation invariance,
and the uniformly bounded half-line restriction on $H^s$,
$s<1/2$, give a coarse $O(\Lambda^s)$ bound for this same tail.
Interpolation gives \eqref{hnew:deep-raw-tail}, the
order-$\nu$ version of \eqref{eq:high-graph-tail}.
The deep rank is $O(\Lambda)$ directly from its selected eigenvalue
bound and \eqref{hnew:trace}.  If its leading asymptotic is desired,
the positive soft-compression argument in
Proposition~\ref{prop:hankel-high-cutoff} applies using
\eqref{hnew:trace}--\eqref{hnew:variance} and the preceding Schur bound;
it gives rank asymptotic to $2c\Lambda/\pi$.
\end{proof}

\begin{lemma}[The order-dependent cross operator]\label{hnew:cross}
Let a raw column $V$ have logarithmic output supported in $t\ge-A$,
with all fixed logarithmic derivatives Hilbert--Schmidt.
For every fixed $j,R$,
\[
 \|(1-p(v/Y))\mathcal H_{c,\nu}V\|_{\mathcal S_2(L^2,H_D^j)}
 \le C_{j,R,A}Y^{-R}\sum_{r\le j+2R}\|\partial_t^rUV\|_2.
\]
Integer $j$ suffice to control the fractional output norms in use.
After endpoint scaling the opposite spectral radius is $Y=\Lambda^2$.
\end{lemma}
\begin{proof}
The Bessel equation gives the exact diagonalization
\begin{equation}\label{hnew:differential}
 \mathcal D_\nu=-\partial_u(u\partial_u)+\frac{\nu^2}{4u},
 \quad \mathcal H_{c,\nu}\mathcal D_\nu=c^2v\mathcal H_{c,\nu},
 \quad
 U\mathcal D_\nu U^{-1}
 =e^{-t}\bigl[-(\partial_t-1/2)^2+\nu^2/4\bigr].
\end{equation}
Its self-adjoint realization is supplied by the angular Fourier
model.  On the tail $v\gg Y$, insert $v^{-R}$ and transfer $v^R$
to $c^{-2R}\mathcal D_\nu^R$.  In logarithmic coordinates this
is $e^{-Rt}$ times a fixed differential polynomial of degree $2R$,
whose coefficients are bounded on the support $t\ge-A$.
Commuting logarithmic derivatives through the smooth tail cutoff
and using dilation reversal proves the estimate.
Approximation by compact cores justifies the diagonalization on
these columns.  Finally
$\mathcal H_{c,\nu}D_\Lambda=D_{1/\Lambda}\mathcal H_{c,\nu}$
turns the dual radius $\Lambda$ into $\Lambda^2$.
This proves the replacement for Lemma~\ref{lem:hi-hankel-cross},
including its polynomial-height form whenever the displayed
derivative budgets are polynomial in height.
\end{proof}

\subsection{Arithmetic inputs and the sharp Mellin passage}

\begin{lemma}[Fixed-form analytic estimates]\label{hnew:arithmetic}
Let $\mu_f$ be the Dirichlet inverse coefficients of $\lambda_f$.
For fixed $f$ there are $b_f>0$ and constants such that
\begin{align}
 |\lambda_f(n)|+|\mu_f(n)|&\le2d(n),\label{hnew:divisor}\\
 \sum_{n\le X}\mu_f(n)&\ll_f Xe^{-b_f\sqrt{\log(2+X)}},\label{hnew:mertens}\\
 |L_f(1/2+i\tau)|&\ll_f(2+|\tau|)^{3/8},
 \qquad \|E_f\psi\|_2\ll_f\|\psi\|_{H_D^{3/8}}.\label{hnew:graph}
\end{align}
For $A_f(X)=\sum_{n\le X}\lambda_f(n)$,
\begin{equation}\label{hnew:mean-square}
 \int_1^X|A_f(y)|^2dy\ll_f X^{3/2},\qquad
 \int_Y^\infty|A_f(y)|^2y^{-2}dy\ll_f Y^{-1/2}.
\end{equation}
The function $L_f$ has polynomial growth on each fixed vertical
strip, a uniform Euler-product lower bound at $2+iT$, and
$O_f(\log(2+|T|))$ zeros in any fixed-radius disk near height $T$.
\end{lemma}
\begin{proof}
For $k\ge2$ and $p\nmid N$, Deligne's primary statement
\citep[Theorem~8.2, p.~302]{deligne1974}
gives unit-modulus normalized Hecke roots.
At $p\mid N$, the newform Euler factor is linear.
The ramified coefficient formulas of
\citet[\S\S1.7--1.8, pp.~509--510]{deligneserre1974}
give $|\lambda_f(p)|\in\{0,1,p^{-1/2}\}$.
For $k=1$, the finite-image representation and the sharper
coefficient bound in
\citet[Theorem~9.1 and equation~(9.3), p.~527]{deligneserre1974}
give $|\lambda_f(n)|\le d(n)$, including the ramified factors.
Thus, in every case, the complete Euler product is
\[
 L_f(s)=\prod_p
     (1-\lambda_f(p)p^{-s}+\xi(p)p^{-2s})^{-1}.
\]
At good primes $|\lambda_f(p^r)|\le r+1$; at bad primes
$\lambda_f(p^r)=\lambda_f(p)^r$ and $|\lambda_f(p)|\le1$.
The inverse coefficients satisfy
$\mu_f(p)=-\lambda_f(p)$, $\mu_f(p^2)=\xi(p)$, and
$\mu_f(p^r)=0$ for $r\ge3$.
In particular $\mu_f(p^2)=0$ at every bad prime.
Multiplicativity proves \eqref{hnew:divisor}, for $f$ and $\bar f$.
The fixed-form Weyl bound of
\citet[Theorem~1.1, p.~2 of the author preprint]{booker2019}
is $O_f(|\tau|^{1/3}\log|\tau|)$, at arbitrary level and
including weight one.  The latter cases are also explicit in
the discussion of dihedral, tetrahedral, octahedral, and
icosahedral forms in footnote~1 on p.~4 of that preprint.
It implies \eqref{hnew:graph} by Plancherel.
Convexity with exponent $1/2+\epsilon$ would not suffice.

Hoffstein and Ramakrishnan's primary theorem
\citep[Theorem~C(3), p.~299]{hoffsteinram1995}
excludes zeros in
\[
 \sigma\ge1-\frac{c_2}{\log(\max(1,|t|)+M_f)}.
\]
For fixed $f$, shrink this to $1-b/\log(3+|t|)$.
On a disk of radius $O(1/\log(3+|t|))$ centered just to the
right of 1, an analytic logarithm exists.
Convexity in that shrinking strip bounds $\log|L_f|$ above by
$O_f(\log\log(3+|t|))$; the Euler logarithm at the center has
that bound in modulus.  Borel--Carath\'eodory on a smaller disk gives
$|L_f(s)^{-1}|\ll_f(\log(3+|t|))^C$ in a narrower region.
Bounded heights follow by compactness, without an exceptional zero.

The divisor hyperbola argument gives
$\sum_{X<n\le X+H}d(n)\ll H\log(2X)+\sqrt X$.
For a half-integer $X$, truncated Perron at $1+1/\log X$ and
height $2\le T\le\sqrt X$ therefore has error
$O(X\log^2X/T+\sqrt X)$: split the near range into
$|n-X|\le X/T$ and its successive doublings.
Moving the contour to $1-b/\log(3+T)$ uses the reciprocal bound.
Choose $T=\exp(d\sqrt{\log X})$ with fixed $d>0$ and absorb the
logarithms into a smaller positive exponent.  This proves
\eqref{hnew:mertens} for half-integers and hence all $X$.
No height-uniform twisted inverse-coefficient assertion is used.

Apply \citet[Theorem~1, p.~62, equation~(57), with setup on
pp.~43--44]{chandrasekharannar1964} to
\[
 \varphi(s)=c^{-s}L_f(s),\quad
 \psi(s)=\varepsilon_fc^{-s}L_{\bar f}(s),\quad
 \Delta(s)=\Gamma(s+a),\quad\delta=1,\quad A_{\rm CN}=1.
\]
The gamma shift, including $a=0$, and complex coefficient
sequences are allowed.  Both Dirichlet scales are $cn$, and
the completed function has the required strip decay.
The inequality $d(n)^2\le d_4(n)$ gives
$\sum_{n\le X}|\lambda_f(n)|^2\ll X\log^3(2X)$.
Their parameters $\beta=1,\beta'=3$ satisfy the included equality
case $2\beta-\delta-1/A_{\rm CN}=0$ and give $\omega_0=1/4$.
After undoing the fixed scale, their error differs from $A_f$
only by the residue constant $L_f(0)$.
The estimate is $O(X^{3/2})+O(X\log^5X)$ and proves the first
part of \eqref{hnew:mean-square}.  Partial integration gives the second.

Finally the functional equation and absolute Euler convergence give
the fixed-strip polynomial bound.  At $2+iT$ both Euler factors
and their reciprocals are bounded uniformly.
Jensen's formula in a larger fixed disk now gives the local zero
count.  Nontrivial zeros lie in $0<\Re s<1$: the Euler product,
the zero-free line at 1, and the functional equation exclude the
other points apart from the usual trivial zeros at $-a-\mathbb N_0$.
\end{proof}

\subsection{A weighted-tent identity and all zero jets}

\begin{proposition}[The required Riesz estimate]\label{hnew:riesz}
For $A_f(X)=\sum_{n\le X}\lambda_f(n)$,
\begin{equation}\label{hnew:riesz-bound}
 R_f(X):=\int_0^XA_f(y)dy=L_f(0)X+O_f(X^{3/4}).
\end{equation}
When $k=1$, $L_f(0)=0$, so the bound is $R_f(X)=O_f(X^{3/4})$.
In particular $A_f(X)=O_{f,\epsilon}(X^{3/8+\epsilon})$ for
every $\epsilon>0$.
\end{proposition}
\begin{proof}
Set
\[
 \phi_X(u)=u^a(X-u)_+,\qquad
 B_a(X)=\sum_{n\le X}\lambda_f(n)n^a(X-n).
\]
The angular lift of $\phi_X$ is $z^\nu(X-|z|^2)_+$ and is smooth
at zero.  Substituting $u=r^2$ and twice using
$\frac d{dr}(r^{\nu+1}J_{\nu+1}(br))
 =br^{\nu+1}J_\nu(br)$ gives Sonine's identity in the required
normalization:
\begin{equation}\label{hnew:sonine}
 (\mathcal H_{c,\nu}\phi_X)(v)
       =\frac{X^{a+1}}{cv}J_{\nu+2}(2c\sqrt{Xv}).
\end{equation}
The outer endpoint is a continuous tent and requires a limiting
argument.  Convolve its compact angular lift with a nonnegative
radial smooth approximate identity on $\mathbb R^2$.
This preserves angular frequency and gives smooth compact angular
profiles, to which \eqref{hnew:poisson} applies.
On the Fourier side it multiplies the transform by a radial
function of modulus at most one, converging pointwise to one.
Equation \eqref{hnew:sonine} and the fixed-order Bessel asymptotic
give the summable majorant $C_{f,X}d(n)n^{-5/4}$ for the dual sum.
Dominated convergence applies there.  The original supports stay
in a fixed compact disk, so their finite theta sums converge,
including the zero tent value at a possible boundary integer.
At logarithmic time zero the identity is therefore
\begin{equation}\label{hnew:weighted-tent}
 B_a(X)=\frac{\varepsilon_fX^{a+1}}c
       \sum_{n\ge1}\frac{\lambda_{\bar f}(n)}n
               J_{\nu+2}(2c\sqrt{Xn})
       =O_f(X^{a+3/4})\quad(X\ge1).
\end{equation}

If $a=0$, then $\nu=0$ and $B_0=R_f$.
Equation \eqref{hnew:weighted-tent} directly proves
$R_f(X)=O_f(X^{3/4})$; no Abel division by a positive weight is
needed.  Completed entireness of
$c^{-s}\Gamma(s)L_f(s)$ implies $L_f(0)=0$.
This proves \eqref{hnew:riesz-bound} in weight one.
The final pointwise argument below applies in this case as well.
For the intervening Abel calculation assume $a>0$.

Twice integrating by parts in the measure with second primitive
$B_a$, against $g(y)=(X-y)y^{-a}$, yields
\begin{equation}\label{hnew:abel-twice}
 \begin{split}
 R_f(X)={}&X^{-a}B_a(X)
  +a(a+1)X\int_0^X B_a(y)y^{-a-2}dy\\
 &+a(1-a)\int_0^X B_a(y)y^{-a-1}dy.
 \end{split}
\end{equation}
There is no boundary at zero because $B_a$ vanishes below 1;
at $X$, $g(X)=0$ and $g'(X)=-X^{-a}$.
The first integral has a finite limit at infinity, since its
integrand is $O(y^{-5/4})$.  Its omitted tail times $X$, and both
other terms, are $O(X^{3/4})$.  Thus
$R_f(X)=CX+O_f(X^{3/4})$.
For $\Re s>1$,
\[
 L_f(s)=s(s+1)\int_1^\infty R_f(y)y^{-s-2}dy
       =C(s+1)+s(s+1)\int_1^\infty(R_f(y)-Cy)y^{-s-2}dy.
\]
The last integral is holomorphic on $\Re s>-1/4$.
Continuing and evaluating at zero gives $C=L_f(0)$.

Finally average $A_f$ on $[X,X+H]$.
Equation \eqref{hnew:riesz-bound} bounds this average by
$O_f(1+X^{3/4}/H)$.  Its difference from $A_f(X)$ is at most
$O_{f,\epsilon}((H+1)X^\epsilon)$, since
$|\lambda_f(n)|\ll_\epsilon n^\epsilon$.
Taking $H=X^{3/8}$ proves the stated pointwise bound.
\end{proof}

For $k\ge2$, the term $L_f(0)X$ cannot in general be discarded.
It is a Riesz residue, not a pole of $L_f$, and creates no extra
theta coordinate.  The direct substitution $J_2\mapsto J_{\nu+2}$
in the order-zero unweighted tent formula would be invalid:
the weight $u^a$ in \eqref{hnew:sonine} is essential.

\begin{proposition}[Sharp Mellin identity and raw tail kernels]
\label{hnew:sharp-mellin}
Let $\phi$ be a selected sharp profile, or an angular Schwartz
profile.  Then for $0<\Re s<1$,
\begin{equation}\label{hnew:improper}
 \lim_{\eta\downarrow0}
 \int_\eta^\infty (E_f\phi)(\log x)x^{s-1/2}\frac{dx}{x}
                =L_f(s)\mathcal M\phi(s).
\end{equation}
The convergence is locally uniform and admits every fixed
$s$-derivative.  For a nontrivial zero $\rho=\beta+i\gamma$ and
$j<m_\rho$, put
\begin{equation}\label{hnew:jet-kernel}
 K_{\rho,j,\Lambda}(u)=
 \left.(\partial_s+L)^j\left[
 u^{s-1}\left(L_f(s)-\sum_{n\le\Lambda u}\lambda_f(n)n^{-s}\right)
                         \right]\right|_{s=\rho}.
\end{equation}
This is a raw $L^2$ kernel annihilating the Hankel band, and
\begin{equation}\label{hnew:jet-exterior}
 \|1_{u>\Lambda}K_{\rho,j,\Lambda}\|_2^2
                           \ll_{f,\rho,j}\Lambda^{-2\beta}.
\end{equation}
\end{proposition}
\begin{proof}
Integration by parts against \eqref{hnew:riesz-bound} proves,
for $\Re s=\beta>0$,
\begin{equation}\label{hnew:integrated-discrepancy}
 \int_X^\infty A_f(y)y^{-s-1}dy
 =\frac{L_f(0)}sX^{-s}+O_{f,s}(X^{-\beta-1/4}).
\end{equation}
The same estimate with fixed logarithmic powers proves local
uniformity of derivatives.
It continues $L_f(s)/s=\int_1^\infty A_f(y)y^{-s-1}dy$
from $\Re s>1$ to $\Re s>0$.
For the disk $1_{(0,b)}$ the theta function is
$\sqrt x A_f(b/x)$, so its omitted lower Mellin integral is
\[
 O_{f,s}(\eta^\beta)
      +O_{f,s}(\eta^{\beta+1/4}b^{-1/4}).
\]
For $b<\eta$, the complete disk integral gives the same majorant.
For a sharp selected profile use
\[
 \phi(u)=\phi(\Lambda-)1_{u<\Lambda}
                      -\int_0^\Lambda\phi'(b)1_{u<b}db.
\]
For $k\ge2$, Proposition~\ref{hnew:geometry} gives
$\phi'(b)=O(b^{a-1})$, where $a\ge1/2>1/4$.
For $k=1$, it gives the stronger $\phi'(b)=O(1)$.
In either case $|\phi'(b)|b^{-1/4}(1+|\log b|)^j$
is integrable at zero.  This proves passage to the limit and
every fixed derivative.  In weight one the $O(\eta^\beta)$
term can be omitted because $L_f(0)=0$; the same bound suffices.
For angular Schwartz profiles use the
same representation over the full positive axis; rapid decay
handles infinity.  This proves \eqref{hnew:improper}.

The kernel \eqref{hnew:jet-kernel} vanishes below $\Lambda^{-1}$,
because the indicated derivatives of $L_f$ vanish at $\rho$.
For $j=0$, Stieltjes summation and the preceding continuation give
\begin{equation}\label{hnew:tail-decomposition}
 K_{\rho,0,\Lambda}(u)
  =-\Lambda^{-\rho}u^{-1}A_f(\Lambda u)
    +\rho u^{\rho-1}\int_{\Lambda u}^\infty
                                A_f(y)y^{-\rho-1}dy.
\end{equation}
The endpoint term is in raw $L^2$ by \eqref{hnew:mean-square}.
Its squared exterior norm is
\[
 \Lambda^{1-2\beta}\int_{\Lambda^2}^\infty
                       |A_f(y)|^2y^{-2}dy\ll_f\Lambda^{-2\beta}.
\]
Equation \eqref{hnew:integrated-discrepancy} makes the second term
$L_f(0)\Lambda^{-\rho}u^{-1}
+O(\Lambda^{-\beta-1/4}u^{-5/4})$.
Its squared exterior norm is
$O(\Lambda^{-2\beta-1}+\Lambda^{-2\beta-2})$.
For $j\ge1$, $\partial_s+L$ cancels the endpoint factor
$\Lambda^{-s}$ and the displayed $L_f(0)$ term.
After the substitution $y=(\Lambda u)z$, the remaining derivatives
insert only fixed powers of $\log z$ in an integrable remainder.
This proves the global $L^2$ assertion and \eqref{hnew:jet-exterior}.

First test the kernels on inverse Hankel transforms of smooth
compact spectral profiles supported strictly inside $(0,\Lambda)$.
They are angular Schwartz profiles.  Equation \eqref{hnew:poisson}
gives their lower theta support gap, and \eqref{hnew:improper}
identifies the lower-tail integral and its jets.
These therefore vanish.  Such profiles are dense in the band;
the raw $L^2$ kernel extends annihilation by continuity.
There is no principal value-at-zero constraint.
\end{proof}

\subsection{Gram control and the deep endpoint inputs}

\begin{proposition}[Full Gram and deep leakage]\label{hnew:gram-deep}
For either selection,
\begin{equation}\label{hnew:gram-bound}
 \|G_\Lambda^{-1}\|\ll_f\Lambda^2\log^2(2+\Lambda).
\end{equation}
For the deep selection,
\begin{equation}\label{hnew:deep-leakage}
 \|(1-S_L)Z_\Lambda^{\mathrm d}\|_{\mathrm{HS}}
                       \ll_f\Lambda^{-17/8}\log(2+\Lambda).
\end{equation}
Every fixed complete window zero packet $\mathcal V_{\mathcal F,L}$
satisfies
\begin{equation}\label{hnew:deep-packets}
 \|Q_\Lambda^{\mathrm d}P_{\mathcal V_{\mathcal F,L}}\|_{\mathrm{HS}}^2
            \ll_{f,\mathcal F}\Lambda^{1-56}(1+L)^C.
\end{equation}
In particular this is smaller than every fixed negative power of $L$.
\end{proposition}
\begin{proof}
The coefficient convolution identity gives finite pointwise inversion
\[
 \psi(u)=\sum_{m\le\Lambda/u}
       \frac{\mu_f(m)}{\sqrt{mu}}(E_f\psi)(\log(mu)).
\]
On $u\ge d/\Lambda$, Minkowski and the change to logarithmic
measure bound its norm by
$\sum_{m\le C\Lambda^2}|\mu_f(m)|m^{-1/2}\|E_f\psi\|_2$.
Use \eqref{hnew:raw-mass-graph} and
$\sum_{m\le Y}d(m)m^{-1/2}\ll\sqrt Y\log(2Y)$.
The result is
$\|\psi\|_2\ll_f\Lambda\log(2+\Lambda)\|E_f\psi\|_2$,
which proves injectivity and \eqref{hnew:gram-bound}.

Time support makes the upper theta tail zero.
Since $\mathcal H_{c,\nu}W_\Lambda\psi=F_\Lambda^{-1}\psi$
is time supported, \eqref{hnew:poisson} makes the lower theta
tail of $W_\Lambda\psi$ zero.  Apply the graph bound to its raw
exterior difference.  Equation \eqref{hnew:deep-raw-tail} with
$\sigma=3/8$, $s=7/16$ has $1-\sigma/s=1/7$.
Summing columns and applying the full Gram bound gives
\[
 \|(1-S_L)Z_\Lambda^{\mathrm d}\|_{\mathrm{HS}}^2
 \ll_f
 \Lambda^{2+1+2(3/8)-56/7}\log^2(2+\Lambda)
 =\Lambda^{-17/4}\log^2(2+\Lambda).
\]
This is the direct substitution in the abstract tail calculation
of Corollary~\ref{cor:hankel-high-theta-tail}, retaining its
additional Gram logarithm.

For a jet in \eqref{hnew:jet-kernel}, band annihilation makes its
pairing with $\psi$ the negative of its exterior pairing with
$(1-P_\Lambda)W_\Lambda\psi$.
The latter has squared raw operator norm $O(\Lambda^{-56})$.
Equations \eqref{hnew:gram-bound} and \eqref{hnew:jet-exterior}
therefore give the squared normalized theta-dual bound
$O(\Lambda^{2-56-2\beta}(1+L)^C)$.
This is a single row functional on the full normalized column;
there is no extra factor equal to the selected rank.
The full Mellin jet vanishes at $\rho$, and the upper theta tail
is zero, so it controls the window jet.
For $\beta<1/2$, left normalization contributes
$\Lambda^{2\beta-1}$; for $\beta>1/2$, right normalization
contributes $\Lambda^{1-2\beta}$ and fixed polynomial factors in $L$.
The respective powers are $\Lambda^{1-56}$ and
$\Lambda^{3-56-4\beta}\le\Lambda^{1-56}$.
Critical polynomial normalization costs only powers of $L$.
The finite packet Gram calculation of
Theorem~\ref{thm:finite-packet}, including repeated zeros and
decoupling of the two endpoints, proves \eqref{hnew:deep-packets}.
\end{proof}

\begin{lemma}[The two one-sided closures]\label{hnew:closure}
The closure of $E_fC_c^\infty((0,1))$ in the negative-half-line
Hardy space is its inner multiple by the pure Blaschke product
of the zeros of $L_f$ in $\Re s>1/2$, with their orders.
The other closure is the reflection by $J$ of the closed
$E_{\bar f}$ core image.  Its model-space divisor is the left
divisor of $L_f$.  No singular inner factor or common exponential
delay is present.
\end{lemma}
\begin{proof}
The core images are square integrable by \eqref{hnew:graph} and
have nonpositive logarithmic support.  Raw contractions give
left translations, so the invariant-subspace theorem applies.
Their Mellin transforms are $L_f(s)\mathcal M\phi(s)$.
At each prescribed point a sufficiently narrow raw bump has
nonzero Mellin transform, so the common zero divisor is exactly
the indicated divisor of $L_f$.
Every generator continues analytically through each finite
boundary point.  An interval where one generator is nonzero
cannot support its singular inner measure; at an isolated
finite-order boundary zero an atom is likewise excluded.
Covering the boundary excludes a common singular measure.
A common exponential inner factor would force all images to
vanish on one fixed interval just below zero.
Take a bump arbitrarily close to raw $1$: when $e^t>1/2$,
only $n=1$ can contribute, and its coefficient is one.
This excludes a common delay.
Finally apply the same proof to $E_{\bar f}$ and reflect its
closure using \eqref{hnew:poisson}.
If $\rho=\beta+i\gamma$ is a left-half zero of $L_f$,
then $1-\rho=1-\beta-i\gamma$ is a right-half zero of
$L_{\bar f}$, with the same order.
The negative-half-line jet
$e^{(1/2-\beta+i\gamma)t}t^j$ for that zero becomes, under $J$,
the positive-half-line jet $e^{(\beta-1/2-i\gamma)t}(-t)^j$.
Thus its Cauchy center is the original $\gamma$, as required.
This uses the dual closure and does not assert that the
finite image projection for $f$ commutes with $J$ or with
$F_\Lambda$.  This proves the holomorphic specialization
of the closure argument in Lemma~\ref{lem:hecke-halfline-closure};
no Hecke-character closure theorem is assumed for $f$.
\end{proof}

\begin{proposition}[Deep local equality and polynomial mass]
\label{hnew:deep-upper}
The deep measure converges vaguely to $\mathfrak b_f$, and
there are fixed $C,A>0$ such that
\begin{equation}\label{hnew:uniform-mass}
 \mu_\Lambda^{\mathrm d}([T,T+1])
             \le C(2+|T|)^A\qquad(\Lambda\ge2,\ T\in\mathbb R).
\end{equation}
It therefore converges on every Schwartz frequency test.
\end{proposition}
\begin{proof}
We verify the hypotheses of
Theorem~\ref{thm:endpoint-local-criterion} explicitly and then
identify which parts of the polynomial-height construction are
unchanged.  Fix $\eta\in C_c^\infty(\mathbb R)$ and set
$F_f(z)=L_f(1/2+iz)$.  Choose a compact frequency interval
containing $\supp\eta$, with endpoints away from critical zeros.
Let $P_{0,L}$ project onto all critical-line window jets in that
interval, with full multiplicities, and put $W_L=S_L-P_{0,L}$.
Give the finitely many real poles of $b=\eta/F_f$ the
negative-time prescriptions $(\tau-\gamma-i0)^{-j}$.
With $T_L=U^{-1}C_bW_L$ one has
\begin{equation}\label{hnew:inverse-column}
 E_fT_L=C_\eta W_L,\qquad
 \|UT_L\|_{\mathcal S_2(L^2,H^j)}\le C_j(1+L)^{d_j}.
\end{equation}
Indeed subtract the compactly cut off principal parts of $b$.
Its inverse transform $\kappa$ has a finite exponential-polynomial
negative tail and a Schwartz remainder, and is Schwartz at
positive infinity.  The input zero moments annihilate the whole
exponential polynomials outside the input interval.
The exterior quadrant kernels are Hilbert--Schmidt, and the
middle kernel has polynomial growth in $L$.
These are the one-dimensional calculations in
\eqref{eq:hi-fixed-bulk}; they involve $F_f$, not the Hankel order.
The exact multiplier identity proves the first assertion of
\eqref{hnew:inverse-column}.

The dual column has these same bounds with $\bar f$: by
\eqref{hnew:mellin-hankel}--\eqref{hnew:poisson},
\begin{equation}\label{hnew:dual-reciprocal}
 \frac{g_{c,\nu}(s)}{L_f(1-s)}
                        =\frac{\varepsilon_{\bar f}}{L_{\bar f}(s)}.
\end{equation}
Reflecting a real-pole prescription and restoring the chosen one
adds only delta derivatives, annihilated by the reflected moments.
Reflection sends the moment conditions precisely to the packet
of the dual reciprocal; both $f$ and $\bar f$ obey
Lemma~\ref{hnew:arithmetic}.  Thus no estimate for a large gamma
quotient is being substituted for this cancellation, and no
Gram-norm symmetry between the two theta images is assumed.

After upper-endpoint translation, the first raw boundary column
has kernel
\[
 (1-p(e^t))\kappa(t-s)1_{[-2L,0]}(s)
\]
before its finite input-jet correction.  Its output is bounded
below and its input nonpositive, so the positive Schwartz tail
of $\kappa$ gives a fixed Hilbert--Schmidt limit, with every
fixed derivative and output polynomial weight.
Lemma~\ref{lem:critical-window-packet} gives strong convergence
to zero of the translated input projections.  Right composition
by these projections therefore removes their corrections in
Hilbert--Schmidt norm.  The dual endpoint follows from
\eqref{hnew:dual-reciprocal}.  Lemma~\ref{hnew:cross} controls both
opposite-endpoint errors with arbitrary powers of $\Lambda^{-1}$.

Here is the exact cutoff algebra.  With $M=M_\Lambda$,
$B=B_\Lambda^p$, and $P_+=1_{t\le L}$,
\begin{align}\label{hnew:cutoff-algebra}
 S_LE_fT_L-E_fMBMT_L
   ={}&S_LE_f(1-M)T_L+S_LE_f(1-B)MT_L\\
      &+P_+E_f(1-M)BMT_L.\nonumber
\end{align}
The term $(P_+-S_L)E_fBMT_L$ is zero by the
band support gap, and $E_fMBMT_L$ has no upper tail.
The cross estimates replace $MT_L$ by $T_L$ in the second
term and $BMT_L$ by $MT_L$ in the last, with a vanishing
Hilbert--Schmidt error after \eqref{hnew:graph}.
Including the two outside-input quadrants of
$S_LC_\eta(1-S_L)$, whose kernel is Schwartz, gives
\begin{equation}\label{hnew:endpoint-decomposition}
 S_LC_\eta(1-P_{0,L})
   =\mathcal F_L+\tau_L\mathcal B_+\tau_{-L}
           +\tau_{-L}\mathcal B_-\tau_L+o_{\mathcal S_2}(1).
\end{equation}
Here $\mathcal B_+$ and $\mathcal B_-$ are fixed Hilbert--Schmidt
operators with ranges in the negative and positive half-lines.
The bulk is $\mathcal F_L=E_fK_\Lambda^pT_L$.
Equation \eqref{hnew:interior}, with
\eqref{hnew:inverse-column}, gives
$\|(1-Q_\Lambda^{\mathrm d})\mathcal F_L\|_2\to0$ directly
in the theta graph norm.  Also $\|\mathcal F_L\|_2$ is polynomial
in $L$: this follows from \eqref{hnew:cutoff-algebra}, its bounded
boundary terms, and the window convolution norm, rather than
from a crude raw bound with a power of $\Lambda$.

For the closure spaces in Lemma~\ref{hnew:closure}, each fixed
core profile, dilated to an endpoint, is approximated in theta
graph norm by the deep selection.  To see this, its angular
Hankel transform is Schwartz; dilation therefore makes the
opposite raw tail arbitrarily small, and
\eqref{hnew:interior} completes the approximation.
Contractivity and density extend it to every fixed vector in
the closure.  Finite exponential-jet packets are dense in the
two Blaschke complements.  Equation \eqref{hnew:deep-packets}
and exponential decay beyond the opposite endpoint give the
required $O_R(L^{-R})$ annihilation on each fixed such packet.
The finite-packet diagonal calculation of
Theorem~\ref{thm:finite-packet} is
\[
 \|P_{K_+}C_\eta\|_2^2+\|P_{K_-}C_\eta\|_2^2
 =\sum_{\beta\ne1/2}m_\rho
 \int|\eta(\tau)|^2
 \frac{|\beta-1/2|\,d\tau}
 {\pi((\tau-\gamma)^2+|\beta-1/2|^2)}.
\]
The local zero count makes these Hilbert--Schmidt norms finite
and justifies exhaustion by finite packets.
Lemma~\ref{lem:critical-window-packet} supplies the critical
contribution $\sum_{\beta=1/2}m_\rho|\eta(\gamma)|^2$.
In particular $C_\eta$ has not been commuted with $P_{0,L}$:
that lemma proves its approximate scalar action on each jet.
All four hypotheses of
Theorem~\ref{thm:endpoint-local-criterion} are now verified.
It gives
\[
 \|(1-Q_\Lambda^{\mathrm d})S_LC_\eta\|_2^2
  =\operatorname{Tr}(D_\Lambda^{\mathrm d}C_\eta C_\eta^*)
       \longrightarrow\int|\eta|^2d\mathfrak b_f.
\]
Polarization, compact mass bounds, and uniform approximation
give vague convergence.

For the uniform estimate, let $H=2+|T|$.
The analytic construction of Lemma~\ref{lem:hi-rectangle}
uses only the fixed-strip bound, Euler-center lower bound,
and $O(\log H)$ local zero count, all proved in
Lemma~\ref{hnew:arithmetic}.  It therefore supplies a bump
$\eta_T=1$ on $[T,T+1]$, a packet of $N_T=O(\log H)$ zeros,
and the proper rational sums of the principal parts of
$z^j/F_f(z)$, with the bounds
\[
 \|\partial_\tau^m(\tau^j\eta_T/F_f-R_{T,j})\|_2
       \le H^{C_{j,m}},\qquad
 \sup_{-A\le t\le0}|R_{T,j,-}(t)-R_{T,j,+}(t)|
       \le H^{C_{j,A}}.
\]
For precision, this construction divides by a disk Blaschke
product and uses Harnack to control the zero-free quotient.
Vertical sides are chosen by averaging logarithmic distances
from zero projections; the horizontal sides have imaginary
parts $\pm3/4$.  Cauchy integrals bound the combined principal
parts, without division by root gaps.  These scalar analytic
steps are independent of the underlying raw transform.

Let $W_{L,T}$ remove the complete packet of window jets, including
off-line zeros.  The one-dimensional division and combined-residue
argument of Lemma~\ref{lem:hi-allpole-boundaries} then gives
\[
 \|U^{-1}C_{\eta_T/F_f}W_{L,T}\|_{\mathcal S_2(L^2,H_D^j)}
       \le H^{C_j}(1+L)^{N_T+C_j},
\]
and polynomially bounded fixed translated endpoint columns.
Each inverse linear factor has a decaying integration kernel
of modulus at most one on a moment-zero input; its cost is
at most $2L$.  At an endpoint, subtract the whole opposite
residue polynomial, which is annihilated exactly.  What remains
is supported in a fixed triangle and is controlled by the
displayed residue-difference bound.
For the dual endpoint use \eqref{hnew:dual-reciprocal}, and for
its raw cross errors use Lemma~\ref{hnew:cross}.
These are the only transform-dependent steps.

Applying \eqref{hnew:interior} to these columns produces a bulk
error of the form
\begin{equation}\label{hnew:height-error}
 H^C(1+L)^{C\log H+C}e^{-\eta_0L},
\end{equation}
with fixed $\eta_0>0$.  The derivative budgets are fixed,
not proportional to $N_T$, and $M=56$ is unchanged.
For $L\ge H^D$, with a fixed sufficiently large $D$, the error
is bounded by one.  Contractivity bounds the boundary residual
by a fixed power of $H$.  The complementary input projection
has rank at most $N_T$ and contributes at most that rank.
Orthogonality of the two input subspaces adds their squared
Hilbert--Schmidt norms exactly.  The distance identity above
therefore gives \eqref{hnew:uniform-mass} for $L\ge H^D$.
For $L<H^D$ use $0\le d\mu_\Lambda^{\mathrm d}/d\tau\le L/\pi$.
Bounded heights follow from the compact limit and the elementary
bound on the remaining bounded cutoff range.
This is the explicit specialization of the argument of
Theorem~\ref{thm:hi-uniform-mass}; no raw strong-convergence
claim replaces an image-distance estimate.
Summing a Schwartz supremum over unit intervals against
\eqref{hnew:uniform-mass} proves Schwartz convergence.
\end{proof}

\subsection{The fixed-threshold functional and the final transfer}

\begin{proposition}[An exterior functional in the theta norm]
\label{hnew:natural-functional}
For the fixed threshold, every $R>0$, and
$q\in L^2((\Lambda,\infty),du)$ extended by zero,
\begin{equation}\label{hnew:exterior-functional}
 \sup_{0\ne\psi\in\operatorname{Ran}\Pi_\Lambda}
  \frac{|\langle W_\Lambda\psi,q\rangle|^2}{\|E_f\psi\|_2^2}
       \le C_{f,R}\Lambda L^{-R}\|q\|_2^2
       \qquad(\supp q\subset(\Lambda,\infty)).
\end{equation}
Equivalently, after relabelling $R$,
\[
 \|(1-P_\Lambda)W_\Lambda j_\Lambda G_\Lambda^{-1/2}\|
                       \le C_{f,R}\Lambda^{1/2}L^{-R}.
\]
\end{proposition}
\begin{proof}
Put $k_q=F_\Lambda^{-1}\Pi_\Lambda P_\Lambda\mathcal H_{c,\nu}q$.
Then $\langle W_\Lambda\psi,q\rangle=\langle\psi,k_q\rangle$ and
$\|k_q\|_2\le\vartheta^{-1/2}\|q\|_2$.
For $J_u(v)=1_{(0,\Lambda)}(v)J_\nu(2c\sqrt{uv})$, evaluation
using $F_\Lambda^{-1}\Pi_\Lambda=F_\Lambda A_\Lambda^{-1}\Pi_\Lambda$
and the exact identity
$P_\Lambda\mathcal H(1-P_\Lambda)\mathcal H P_\Lambda=P_\Lambda-A_\Lambda$
give
\begin{align}\label{hnew:row-energy}
 |k_q(u)|^2
 &\le c^2\|q\|_2^2
 \langle J_u,A_\Lambda^{-2}\Pi_\Lambda(1-A_\Lambda)J_u\rangle\\
 &\le\vartheta^{-2}\|q\|_2^2
                 \int_\Lambda^\infty|B_\Lambda(v,u)|^2dv.\nonumber
\end{align}
Commutation here is only that of spectral functions of $A_\Lambda$.
Replacing $J_u$ by $u\partial_uJ_u$ gives the differentiated bound.
The raw Lommel kernel is
\[
 B_\Lambda(v,u)=\frac{c\sqrt\Lambda}{v-u}
 \left[\sqrt v J_{\nu+1}(2c\sqrt{\Lambda v})J_\nu(2c\sqrt{\Lambda u})
 -\sqrt u J_{\nu+1}(2c\sqrt{\Lambda u})J_\nu(2c\sqrt{\Lambda v})\right].
\]
For $v>\Lambda$, $u\le\Lambda/2$, its modulus is at most
$C\Lambda^{1/4}v^{-3/4}(1+\Lambda u)^{-1/4}$.
Using $zJ_\nu'(z)=\nu J_\nu(z)-zJ_{\nu+1}(z)$ changes the last
exponent to $1/4$ after $u\partial_u$.
The integral of $\Lambda^{1/2}v^{-3/2}$ over $v>\Lambda$ is bounded.
Thus \eqref{hnew:row-energy} proves
\begin{equation}\label{hnew:row-bounds}
 |k_q(u)|\ll\|q\|_2(1+\Lambda u)^{-1/4},\qquad
 |u k_q'(u)|\ll\|q\|_2(1+\Lambda u)^{1/4}.
\end{equation}
These estimates, including the exact energy step, replace
Lemma~\ref{nh:row}; a bounded raw inverse alone would not imply them.

Finite coefficient inversion gives the theta-dual representative
\[
 \beta_k(v)=\sqrt v\sum_{m\ge1}\frac{\overline{\mu_f(m)}}m k(v/m),
       \quad 0<v<\Lambda,
\]
extended by zero above $\Lambda$.
Separate $m=1$, whose squared $L^2(dv/v)$ norm is $\|k\|_2^2$.
Put $X=\Lambda v$, and for large $X$ split the remaining terms
at $Y=X/(\log X)^D$.  The small part is bounded by
\[
 C\|q\|_2X^{-1/4}\sum_{m\le Y}d(m)m^{-3/4}
       \ll\|q\|_2(\log X)^{1-D/4}.
\]
The extra logarithm is retained and absorbed by choosing $D$
sufficiently large.  For the rest use \eqref{hnew:mertens} and
Abel summation with
\[
 \left|\frac d{dy}\{y^{-1}k(v/y)\}\right|
       \ll\|q\|_2y^{-2}(1+X/y)^{1/4}.
\]
The resulting integral is bounded by
$C\|q\|_2\int_Y^\infty e^{-b_f\sqrt{\log(2+y)}}
(1+X/y)^{1/4}dy/y$, and is smaller than every fixed
negative power of $\log X$.
For bounded $X$ it is bounded.
Consequently
$|\beta_{k,\ge2}(v)|\ll_R\|q\|_2\sqrt v
[\log(2+\Lambda v)]^{-R}$.
Integrating first over $v\le\sqrt\Lambda$ and then over its
complement gives
\[
 \|\beta_k\|_{L^2(dv/v)}^2\ll_R\Lambda L^{-R}\|q\|_2^2.
\]

To justify the dual identity, first integrate
$\langle\psi,k\rangle$ over $u\ge\epsilon$.
Finite inversion and $v=mu$ give the same expression with the
sum in $\beta_k$ truncated at $m\le v/\epsilon$.
It converges uniformly on each $v\ge\delta>0$ by Abel summation.
Near zero, if $a>0$, angular regularity gives
$k(v/m)=O_\Lambda((v/m)^a)$, so the truncated representatives
are bounded by $C_\Lambda v^{a+1/2}\sum_m d(m)m^{-1-a}$.
For $a=0$ this absolute series would diverge and is not used.
Instead $k(v/m)=k(0)+O_\Lambda(v/m)$.
Abel summation in \eqref{hnew:mertens} gives
\[
 \sup_{M\ge1}\left|\sum_{m\le M}\frac{\mu_f(m)}m\right|<\infty,
 \qquad \sum_{m\ge1}\frac{d(m)}{m^2}<\infty.
\]
The truncated representatives are therefore
$O_\Lambda(\sqrt v)$ in weight one, uniformly in their
truncation.  Both bounds give integrable squared
$L^2(dv/v)$ majorants near zero.
Thus they converge in that space and
\[
 \langle\psi,k\rangle
   =\langle (E_f\psi)(\log v),\beta_k(v)\rangle_{L^2(dv/v)}.
\]
Cauchy--Schwarz proves \eqref{hnew:exterior-functional}.
\end{proof}

\begin{proposition}[Natural packets and smoothed leakage]
\label{hnew:natural-transfer-inputs}
For every fixed complete zero packet and every $R$,
\begin{equation}\label{hnew:natural-packets}
 \|Q_\Lambda^{(\vartheta)}P_{\mathcal V_{\mathcal F,L}}\|_{\mathrm{HS}}^2
                                             =O_{f,\mathcal F,R}(L^{-R}).
\end{equation}
For $g\in C_c^\infty(\mathbb R)+\mathbb C(\tfrac12e^{-|t|})$,
\begin{equation}\label{hnew:smoothed-leakage}
 \|V(g)(1-S_L)Z_\Lambda^{(\vartheta)}\|_{\mathrm{HS}}
                                                =O_{f,g,R}(L^{-R}).
\end{equation}
For every $h\in C_c^\infty(\mathbb R)$,
\begin{equation}\label{hnew:trace-transfer}
 \operatorname{Tr}((Q_\Lambda-S_LQ_\Lambda S_L)V(h))\longrightarrow0.
\end{equation}
\end{proposition}
\begin{proof}
For a jet kernel from \eqref{hnew:jet-kernel}, annihilation of the
band replaces its pairing with $\psi$ by its exterior pairing
with $W_\Lambda\psi$.
Equations \eqref{hnew:exterior-functional} and
\eqref{hnew:jet-exterior} give the squared theta-dual bound
$O(\Lambda^{1-2\beta}L^{-R})$.
The full Mellin jet vanishes and the upper theta tail is zero.
Left-endpoint normalization contributes $\Lambda^{2\beta-1}$;
right-endpoint normalization contributes $\Lambda^{1-2\beta}$,
giving $\Lambda^{2-4\beta}\le1$ for $\beta>1/2$.
Central jets cost only powers of $L$.
The finite Gram calculation of Theorem~\ref{thm:finite-packet}
therefore proves \eqref{hnew:natural-packets}.

For completeness the transform-independent raw kernel used in
Lemma~\ref{nh:raw-tail} is, in output coordinate $r=t+L$,
\[
 K_g(r,u)=u^{-1/2}\int_{\Lambda u}^\infty
   y^{-1/2}g(r+\log(y/(\Lambda u)))\,dA_f(y),\qquad u>\Lambda.
\]
Integration by parts gives an endpoint
$-\Lambda^{-1/2}u^{-1}A_f(\Lambda u)g(r)$ and the integral
\[
 -\Lambda^{-1/2}u^{-1}\int_1^\infty A_f(\Lambda uz)z^{-3/2}
                    (g'-g/2)(r+\log z)dz.
\]
Equation \eqref{hnew:mean-square} bounds their combined
Hilbert--Schmidt norm by
\begin{equation}\label{hnew:raw-smoothed-tail}
 \|V(g)1_{t<-L}E_f(1-P_\Lambda)\|_{\mathrm{HS}}
       \le C_f\Lambda^{-1/2}(\|g\|_2+\|g'\|_2).
\end{equation}
The Minkowski integral is a constant times
$\int_1^\infty z^{-5/4}dz$.
First use compact smooth $g$, then approximate in $L^1\cap H^1$.
Approximation on the raw graph core identifies the kernel
extension with the displayed theta operator.

The upper tail of $E_f\psi$ and the lower tail of $E_fW_\Lambda\psi$
are zero.  Thus the normalized lower tail in
\eqref{hnew:smoothed-leakage} is the composition of
\eqref{hnew:raw-smoothed-tail} with
$(1-P_\Lambda)W_\Lambda j_\Lambda G_\Lambda^{-1/2}$.
Its factors $\Lambda^{-1/2}$ and $\Lambda^{1/2}$ cancel, proving
the asserted arbitrary logarithmic rate.
The Green kernel $\kappa(t)=\tfrac12e^{-|t|}$ lies in
$L^1\cap H^1$, so it is covered.

The purely operator-theoretic proof of
Proposition~\ref{nh:trace-transfer} now applies with these verified
leakage hypotheses.  Explicitly, put $b=(1-\partial_t^2)h$, so
$h=b*\kappa$, and $P_o=1-S_L$.
The isometry of $Z_\Lambda$ gives
$\|V(g)S_LZ_\Lambda\|_{\mathrm{HS}}\le\sqrt{2L}\|g\|_2$.
Cyclicity and $Q-S_LQS_L=P_oQ+S_LQP_o$ yield
\begin{align*}
 |\operatorname{Tr}((Q-S_LQS_L)V(h))|
 &\le\|V(b^*)Z_\Lambda\|_2\|V(\kappa)P_oZ_\Lambda\|_2\\
 &\quad+\|V(b^*)P_oZ_\Lambda\|_2
                         \|V(\kappa)S_LZ_\Lambda\|_2\\
 &\ll_h L^{1/2-R}+L^{-2R}.
\end{align*}
The norms in this last display are Hilbert--Schmidt norms.
This proves \eqref{hnew:trace-transfer} without a rank loss.
\end{proof}

\begin{proof}[Proof of Theorem~\ref{hnew:theorem}]
The rank and Gram assertions are
Propositions~\ref{hnew:geometry} and~\ref{hnew:gram-deep}.
Proposition~\ref{hnew:deep-upper} gives the full deep measure limit
and its uniform polynomial mass bound.  Its uncompressed trace
follows already from
\[
 \|Q_\Lambda^{\mathrm d}-S_LQ_\Lambda^{\mathrm d}S_L\|_1
 \le2\sqrt{\operatorname{rank}Q_\Lambda^{\mathrm d}}\,
       \|(1-S_L)Z_\Lambda^{\mathrm d}\|_2
       =O_f(\Lambda^{-13/8}\log(2+\Lambda)).
\]

For the fixed threshold and a fixed finite zero packet, put
$P_V=P_{\mathcal V_{\mathcal F,L}}$.
Positivity, $0\le D_\Lambda\le S_L$, and
\eqref{hnew:natural-packets} give
\[
 \|(1-D_\Lambda^{1/2})P_V\|_{\mathrm{HS}}^2
       \le\operatorname{Tr}(P_V(1-D_\Lambda)P_V)
       =\|Q_\Lambda P_V\|_{\mathrm{HS}}^2=o(1).
\]
For a nonnegative compact frequency test $\varphi$, it follows that
\[
 \operatorname{Tr}(D_\Lambda C_\varphi)
 \ge\|P_VD_\Lambda^{1/2}C_\varphi^{1/2}\|_{\mathrm{HS}}^2
 =\operatorname{Tr}(P_VC_\varphi)+o(1).
\]
Theorem~\ref{thm:finite-packet} identifies the last limit.
Increasing the finite packet gives the all-zero lower bound;
the local zero count gives the necessary Cauchy-tail summability.
For large $\Lambda$ the deep raw space is contained in the
fixed-threshold space, hence their theta-image projections are
nested and
\[
 0\le D_\Lambda^{(\vartheta)}\le D_\Lambda^{\mathrm d}.
\]
The deep local limit gives the matching upper bound and its
polynomial local mass passes to the fixed threshold.
Compact convergence and this mass bound imply Schwartz convergence.
This is the positive sandwich and exhaustion argument of
Theorem~\ref{nh:theorem}, now with its lower-bound and upper-bound
inputs explicitly established for $f$.
The compressed Fourier trace identity and
\eqref{hnew:trace-transfer} prove \eqref{hnew:trace-theorem}.

The weighted zero count
$\sum_\rho m_\rho/(1+\gamma^2)<\infty$ follows from the local
zero bound.  The functional equation and
$L_{\bar f}(s)=\overline{L_f(\overline s)}$
reflect $\rho$ to $1-\overline\rho$ in the same $L_f$ divisor,
with the same order and ordinate.  Reality of the coefficients
is not assumed.
The assumed zero-free half-plane therefore confines the divisor
to $|\beta-1/2|\le\delta$.
Theorem~\ref{thm:band-bound}, whose hypothesis is precisely this
band and weighted count, proves the final inequality.
This last deduction changes neither its constants nor its
test-function norm.
\end{proof}

\begin{corollary}[Odd two-dimensional Artin functions]
\label{hnew:odd-artin}
Let $\rho:G_{\mathbb Q}\to GL_2(\mathbb C)$ be a continuous
irreducible odd Artin representation, with Artin conductor $N$.
Let $\lambda_\rho(n)$ be the Dirichlet coefficients of its
complete finite-prime Artin Euler product.
Use the raw Hankel transform
$\mathcal H_{2\pi/\sqrt N,0}$ and the theta operator
\[
 (E_\rho\psi)(t)
       =e^{t/2}\sum_{n\ge1}\lambda_\rho(n)\psi(ne^t).
\]
Then all the fixed-threshold and deep measure, trace, uniform-mass,
and half-plane assertions of Theorem~\ref{hnew:theorem} hold
for $L(s,\rho)$, with leading term $2\log\Lambda$ and no polar
coordinates.  This includes nonsolvable projective image $A_5$.
\end{corollary}
\begin{proof}
The classical modularity theorem of
\citet[Theorem~10.1(ii) and Corollary~10.2(ii), author-preprint
pp.~20--21]{kharewintenberger2009}
associates $\rho$ to a primitive holomorphic newform $f$ of
weight one.  By
\citet[Theorem~4.6, p.~514; proof pp.~515--516]{deligneserre1974},
its level is exactly the Artin conductor and
\begin{equation}\label{hnew:artin-all-primes}
 L(s,\rho)=L_f(s)
 =\prod_p\det(1-\rho(\operatorname{Frob}_p)p^{-s}
                       \mid V_\rho^{I_p})^{-1},
\end{equation}
including every ramified prime.
The representations are identified by their good Frobenius traces
and semisimplicity; this permits the full conductor and local-factor
statement, not merely an equality after deleting bad primes.
The same source's Corollary~4.8 gives entireness.

Oddness gives eigenvalues $1,-1$ at complex conjugation, so the
archimedean factor is
$\Gamma_{\mathbb R}(s)\Gamma_{\mathbb R}(s+1)
=2(2\pi)^{-s}\Gamma(s)$.
It matches the weight-one normalization up to the irrelevant
constant 2, and fixes $c=2\pi/\sqrt N$, $\nu=0$.
Equality \eqref{hnew:artin-all-primes} identifies every Dirichlet
coefficient.  Therefore the raw concentration operator, selected
space, theta map, full Gram matrix, and image projection are
literally those of $f$.  Theorem~\ref{hnew:theorem} applies.
The dual representation supplies the opposite boundary through
the functional equation; it is not added as a second theta channel.
The Artin--modular identification is classical.  The corollary
identifies its finite-cutoff realization.
\end{proof}

The theorem fixes the weight, level, and nebentypus.  Its constants
need not be uniform as the form varies.  The deep rank is asymptotic
to $4\Lambda/\sqrt N$; the logarithmic rank remainder applies to a
fixed threshold.  Even two-dimensional Artin representations and
general higher-degree functions require further cutoff estimates.
